% Tapis de Quillen, with Quillen's offprint and preprint (cote 162-5, typed part): archive cote 162-5, site page 54 (PDF page 55).
% offprint. Draft compared: none.
% AI-transcribed by Claude Opus 5.5 (claude-opus-5-5, adaptive thinking, effort high), Message Batches API, batch t162_full (msgbatch_01UAh3zyWPd454AeDkhnmQ33), 2026-11-13; not yet checked by eye.
\scanpage{162-5/54}
\mspage{On the Representation of Hermitian Forms as Sums of Squares \quad 241}

Corollary 1. Let
\[ E^{\wedge}(\zeta)=\sum_{|\alpha|=|\beta|=m} a_{\alpha\beta}\,\zeta^{\alpha}\,\zeta^{*\beta} \]
be a hermitian form homogeneous of degree $m$ on $\mathbf{C}^n$ which is strictly
positive, i.e. $E^{\wedge}(\zeta)>0$ if $\zeta\neq 0$. Then for sufficiently large $N$,
\[ E^{\wedge}(\zeta)=\frac{\sum\limits_{i}|P_i(\zeta)|^2}{\sum\limits_{|\gamma|=N}(\gamma!)^{-1}\,|\zeta^{\gamma}|^2} \]
where $P_i(\zeta)$ is a basis for the holomorphic polynomials of degree $m+N$ on $\mathbf{C}^n$.

Corollary 2. Let
\[ E^{\sim}(f)=\sum_{|\alpha|=|\beta|=m} a_{\alpha\beta}\,D^{\alpha}f(D^{\beta}f)^* \]
be a differential quadratic form in $n$ variables homogeneous of degree
$(m,m)$, such that $a_{\beta\alpha}=a^*_{\alpha\beta}$. If $E^{\sim}$ has no non-zero characteristics, i.e.
$\sum a_{\alpha\beta}\,\zeta^{\alpha}\,\zeta^{*\beta}>0$ for $\zeta\neq 0$, then for sufficiently large $N$ we may write
\[ \sum_{|\gamma|=N}(\gamma!)^{-1}\,E^{\sim}(D^{\gamma}f)=\sum_{i}|P_i(D)\,f|^2 \]
where $P_i(D)$ is a basis for the homogeneous polynomials of degree $m+N$
in $D_1,\dots,D_n$.

These corollaries are just reinterpretations of the implication (i)$\Rightarrow$(ii)
of the theorem.

Example. Let $E^{\sim}(f)=|D_1^2\,u|^2-|D_1\,D_2\,u|^2+|D_2^2\,u|^2$ where $n=2$. The dif-%
ferential quadratic form $E^{\sim}$ is not a sum of squares because $E^{\sim}(z^1\,z^2)<0$.
But $E^{\sim}(\zeta)=(|\zeta_1|^2-|\zeta_2|^2)^2+|\zeta_1\,\zeta_2|^2>0$ if $\zeta\neq 0$ so the theorem applies.
In fact $E^{1\sim}(f)=|D_1^3\,f|^2+|D_2^3\,f|^2$ and $E^N$ is positive definite for $N\geqq 3$.

Corollary 3. (Projective Nullstellensatz over $\mathbf{C}$.) If $Q_j(z)$ $1\leqq j\leqq l$ are
homogeneous polynomials in $z^1,\dots,z^n$ without a common zero $\zeta\neq 0$, then
the ideal generated by the $Q_j$ contains all homogeneous polynomials of
sufficiently high degree.

Proof. We may assume that the $Q_j$ have the same degree $m$. Let
\[ E^{\wedge}(\zeta)=\sum_{j}|Q_j(\zeta)|^2; \]
then $E^{\wedge}(\zeta)>0$ if $\zeta\neq 0$ so by Corollary 2 for all sufficiently high $N$, there
is an expression
\[ \sum_{j}\sum_{|\gamma|=N}(\gamma!)^{-1}\,|D^{\gamma}\,Q_j(D)\,f|^2=\sum|P_i(D)\,f|^2 \]
where $P_i(D)$ is a basis for the space $W$ of homogeneous polynomials
of degree $m+N$. It follows that the set $D^{\gamma}\,Q_j(D)$ for $|\gamma|=N$, $1\leqq j\leqq l$
spans $W$ for otherwise we could take $f$ to be a non-zero polynomial in

\mspage{\hand{54}}
